\section{Estado atual do aplicativo}
\label{sec:inventario}

Esta seção corresponde à primeira etapa do trabalho (Seção~\ref{sec:etapas}). A quantidade de campos de cada
formulário está na Tabela~\ref{tab:formularios}.

\subsection{Arquitetura e armazenamento}
\label{sec:inv-arquitetura}

O aplicativo tem três partes. A interface, em Flutter, fica em \arq{lib/}: os 15 arquivos de tela estão em
\arq{lib/Presentation/Screens/}, e os modelos, os \textit{providers} e os serviços de acesso à API, em
\arq{lib/Core/}. A API, em FastAPI, concentra as rotas em \arq{Backend/main.py}, e as rotas da CIF ficam em um
módulo próprio, \arq{Backend/cif_api.py}, incluído na linha 37. O banco PostgreSQL é criado por
\arq{Backend/schema.sql}, com 14 tabelas, e a tabela da CIF vem de uma migração separada,
\arq{Backend/migrations/001_create_avaliacoes_cif.sql}.

Cada um dos outros seis formulários envia a avaliação preenchida em uma única requisição \texttt{POST}, e a API grava
o registro. Não há rota para ler essas avaliações de volta: fora da CIF, as únicas rotas de consulta são a raiz da
API e a lista de pacientes (\arq{Backend/main.py}, linhas 39 e 549). Só a CIF tem rotas para listar as avaliações de
um paciente e consultar uma delas (\arq{Backend/cif_api.py}, linhas 405 e 421). Por isso, histórico e comparação
existem apenas na CIF. O Quadro~\ref{qua:inv-formularios} mostra as tabelas de cada formulário.

\begin{quadro}[htb]
	\centering
	\caption{Tabelas do banco de cada formulário}
	\label{qua:inv-formularios}
	\footnotesize
	\begin{tabularx}{\textwidth}{|>{\raggedright\arraybackslash}p{0.22\textwidth}|>{\raggedright\arraybackslash}X|}
		\hline
		Formulário & Tabelas \\ \hline
		CIF & \arq{avaliacoes_cif} \\ \hline
		ASIA & \arq{avaliacoes_asia}, \arq{scores_motores}, \arq{scores_sensoriais} \\ \hline
		GAS & \arq{avaliacoes_gas}, \arq{metas_gas} \\ \hline
		MEEM & \arq{avaliacoes_meem} \\ \hline
		Eletrodiagnóstico & \arq{avaliacoes_eletrodiagnostico}, \arq{medicoes_musculares} \\ \hline
		Densitometria & \arq{avaliacoes_densitometria}, \arq{densitometria_regioes}, \arq{densitometria_tendencias} \\ \hline
		Anamnese & \arq{anamneses} \\ \hline
	\end{tabularx}
	\fonte{Elaborado pelo autor a partir de \arq{Backend/schema.sql} e da migração da CIF \cite{higa2026bemtevi}. As
	tabelas \arq{pacientes} e \arq{usuarios} completam as 14 do esquema.}
\end{quadro}

Cada formulário tem uma tabela principal, ligada ao paciente por \arq{paciente_id}, e os blocos repetidos ficam em
tabelas filhas: os miótomos e dermátomos da ASIA, as metas da GAS, os grupos musculares do eletrodiagnóstico e as
regiões e tendências da densitometria. A tabela da CIF guarda o estado da avaliação, rascunho ou concluída, as
versões do catálogo e das regras e as respostas e os resultados em JSON
(\arq{001_create_avaliacoes_cif.sql}, linhas 4--17). As tabelas dos outros seis formulários não têm estado, versão
nem data de registro. Nenhuma tabela guarda o identificador de quem aplicou a avaliação. ASIA, eletrodiagnóstico e
MEEM gravam o nome do examinador como texto (\arq{schema.sql}, linhas 28, 123 e 146), e GAS, anamnese e
densitometria não registram o avaliador. O esquema também deixa espaço para dados duplicados ou inconsistentes. O
CPF não tem restrição de unicidade (linha 7), e a anamnese repete nome, celular, e-mail e data de nascimento que já
estão no cadastro do paciente (linhas 111--114). As medidas do eletrodiagnóstico são numéricas (linhas 131--136),
mas o índice de acomodação é gravado como texto (linhas 138--139). Quando falta a reobase ou a acomodação, o aplicativo envia o
texto ``N/A'' no lugar do índice (\arq{eletroFormProvider.dart}, linhas 29 e 38), e a ausência do dado fica
registrada como se fosse um valor.

\subsection{Problemas de dados}
\label{sec:inv-problemas}

Os problemas encontrados recebem um identificador, de PD1 a PD5, que a Seção~\ref{sec:especificacao} usa como fonte
dos requisitos. Os quatro primeiros já foram mencionados na Seção~\ref{sec:problema}. O Quadro~\ref{qua:inv-problemas}
resume os cinco, com a localização no código, a dimensão de qualidade afetada (Seção~\ref{sec:ref-qualidade}) e os
requisitos propostos para corrigi-los. O PD1 e o PD5 são detalhados em seguida.

\begin{quadro}[H]
	\centering
	\caption{Problemas de dados encontrados no código do aplicativo}
	\label{qua:inv-problemas}
	\footnotesize
	\begin{tabularx}{\textwidth}{|l|>{\raggedright\arraybackslash}X|>{\raggedright\arraybackslash}p{0.24\textwidth}|>{\raggedright\arraybackslash}p{0.13\textwidth}|>{\raggedright\arraybackslash}p{0.11\textwidth}|}
		\hline
		& Problema & Código (arquivo, linhas) & Dimensão & Requisitos \\ \hline
		PD1 & Paciente identificado pelo nome. Sem correspondência exata, a API cria um cadastro só com o nome &
			\arq{main.py}, 54--61, 134--140, 220--227, 271--278, 344--351 e 675--682 & Identificação do paciente &
			RF08, RN12 \\ \hline
		PD2 & Comentários da ASIA descartados. O \textit{provider} guarda o texto, mas o objeto enviado não o inclui, e
			nem o modelo da API nem a tabela \arq{avaliacoes_asia} têm onde guardá-lo & \arq{asiaForm.dart}, 410--417;
			\arq{asiaFormProvider.dart}, 22 e 205; \arq{models.py}, 15--22 & Completude & RF11, RN14 \\ \hline
		PD3 & Avaliador fixo no MEEM. O examinador é enviado sempre como ``Dr. Examinador'', qualquer que seja o
			usuário que registrou a avaliação & \arq{meemForm.dart}, 82 & Correção & RF09, RN15 \\ \hline
		PD4 & Data e paciente fictícios na densitometria. A data gravada como data do exame é a do envio, e, sem
			paciente selecionado, o nome enviado é ``Paciente Teste'', que a regra do PD1 transforma em cadastro &
			\arq{dex.dart}, 133 e 134 & Correção & RF12, RN12, RN15 \\ \hline
		PD5 & Escore da GAS calculado de duas formas. A tela e o servidor usam a mesma fórmula, mas o servidor troca por
			zero os valores ponderados negativos da linha de base e do resultado alcançado, e a tela não faz essa troca & \arq{gasCalculator.py}, 16--20;
			\arq{GoalsModels.dart}, 48--49 & Concordância & RN19, RNF05 \\ \hline
	\end{tabularx}
	\fonte{Elaborado pelo autor a partir de \citeonline{higa2026bemtevi}. Dimensões de
	\citeonline{weiskopf2013dataquality}; o PD1 é um problema de identificação do paciente, do tipo que gera cadastros
	duplicados \cite{just2016patientmatching}.}
\end{quadro}

\textbf{PD1. Paciente identificado pelo nome.} Os outros seis formulários enviam o nome do paciente, e não seu
identificador. A API pesquisa o nome por correspondência exata e, se não o encontrar, cria um cadastro apenas com o
nome. Basta um erro de digitação, uma letra maiúscula ou um espaço a mais para que seja criado um paciente diferente
e o histórico seja dividido. As telas preenchem o nome com base no paciente selecionado, como em \arq{meemForm.dart},
linha 251, e em \arq{eletroForm.dart}, linha 269, mas o campo permanece editável. O identificador do paciente
selecionado está disponível no aplicativo (\arq{patientProvider.dart}, linha 10), e a CIF o utiliza. Na GAS, o nome
do paciente é inserido no campo ``Plano de intervenção do cliente'' (\arq{gasForm.dart}, linhas 51 e 407--412), e o
conteúdo desse campo é enviado simultaneamente como identificador do paciente e como plano de intervenção (linhas 109
e 112). Se o usuário digitar o plano nesse campo, conforme indicado pelo rótulo, a API cria um paciente cujo nome é o
texto do plano. A rota que atualiza o cadastro também pesquisa o paciente pelo nome (\arq{Backend/main.py}, linhas
605--613).

\textbf{PD5. Escore da GAS calculado de duas formas.} Na escala usada pelo aplicativo, que vai de $-2$ a $+2$, o
nível $-1$ é a linha de base (\arq{gasForm.dart}, linhas 1140--1157), de modo que valores ponderados negativos são o
caso comum. Com uma meta fictícia de importância 2, dificuldade 2, linha de base $-1$ e resultado alcançado 0, o
servidor grava escore inicial 50, escore alcançado 50 e evolução 0. A conta da tela dá 40, 50 e $+10$
(\arq{entregas/scripts/conferir_calculo_gas.py}). Depois do envio, a tela substitui os seus valores pelos do servidor
(\arq{gasForm.dart}, linhas 117--126). Um paciente que sai da linha de base e atinge o resultado esperado fica, no
banco, sem evolução.

